\documentclass[acmlarge,nonacm]{acmart}

\settopmatter{printacmref=false}
\renewcommand\footnotetextcopyrightpermission[1]{}
\pagestyle{plain}
\setcopyright{none}

\usepackage{enumitem}
\usepackage{multicol}
\setlist{nosep,leftmargin=1.5em}
\setlength{\parskip}{2pt}
\setcounter{tocdepth}{1}

\begin{document}

\title{Procedural Noise and Textures}
\subtitle{Computer Graphics Course Project: Summary and Plan (Team 11, CG\_04)}

\author{Keshav Singh}
\email{keshav24300@iiitd.ac.in}
\affiliation{\institution{IIIT Delhi}\country{India}}

\author{Karanveer Singh}
\email{karanveer24294@iiitd.ac.in}
\affiliation{\institution{IIIT Delhi}\country{India}}

\begin{abstract}
We will implement gradient noise (with value and simplex noise as alternative bases) and use it to build fBm, turbulence, marble, wood, cloud and terrain on the GPU, add noise-based bump mapping, and study how noise type, octaves and persistence affect quality and cost.
\end{abstract}

\maketitle

\newpage{}
\begingroup
\small
\renewcommand\contentsname{Index}
% \begin{multicols}{2}
\tableofcontents
% \end{multicols}
\endgroup
\newpage{}

\section{Project Topic and Reference Paper}\label{sec:topic}
\textbf{Procedural Noise and Textures (CG\_04).} Reference: K.~Perlin, ``An Image Synthesizer,'' \emph{Proc.\ ACM SIGGRAPH '85}, Computer Graphics 19(3), pp.~287--296, DOI 10.1145/325334.325247~\cite{perlin1985}. We also use Perlin's ``Improving Noise''~\cite{perlin2002} for the revised interpolant and Inigo Quilez's articles~\cite{iq} for practical fBm, domain-warping and terrain techniques.

\section{Problem and Objective}\label{sec:problem}
Image-based textures need storage, tile visibly and are hard to vary. The paper instead synthesizes surface detail procedurally: a band-limited pseudo-random \emph{gradient noise} $n(\mathbf{x})$ is evaluated per sample and scaled copies are summed to mimic the self-similar structure of natural materials, giving resolution-independent, parameterizable textures at negligible memory cost. We will implement this pipeline in a real-time shader framework, produce a small family of materials, and study quantitatively how the noise parameters control appearance and cost.

\section{Primary Implementation}\label{sec:primary}
\begin{itemize}
\item \textbf{Gradient noise} in 2D and 3D (lattice gradients, dot products with offsets, smooth interpolation), with both the original cubic $3t^2-2t^3$ and the quintic $6t^5-15t^4+10t^3$ fade.
\item \textbf{Value noise} as a baseline behind the same interface, so any base noise can feed the fractal sums.
\item \textbf{Fractal sums:} fBm $\sum_i p^i n(2^i\mathbf{x})$ with adjustable octaves and persistence $p$, and turbulence $\sum_i p^i |n(2^i\mathbf{x})|$.
\item \textbf{Textures:} marble (sine plus turbulence), wood (perturbed rings) and clouds (fBm with a colour ramp) on 3D objects using object-space coordinates.
\item \textbf{Basic terrain:} an fBm height field on a displaced grid mesh, finite-difference normals, height/slope colouring (water, grass, rock, snow) and directional lighting.
\end{itemize}

\section{Proposed Extension}\label{sec:ext}
\textbf{Main: noise-driven bump mapping with animation.} We perturb the surface normal using the gradient of the fBm height field (analytic or finite-difference) before Phong lighting, giving apparent relief without extra geometry; time as a third noise axis animates clouds and surfaces.
\textbf{Alternatives} (chosen by progress): (a)~\emph{domain warping}, $f(\mathbf{p}+\mathbf{w}(\mathbf{p}))$ with $\mathbf{w}$ an fBm, for flowing marble, smoke and terrain; (b)~2D/3D simplex noise as a further base; (c)~erosion-like derivative-based terrain.

\section{Proposed Experiment}\label{sec:exp}
\textbf{Question:} how do octave count $N$, persistence $p$ and the base noise type (value, Perlin, simplex) affect visual quality and cost, and beyond which octave does extra detail stop being perceptible?

\textbf{Varied:} $N\in\{1,\dots,12\}$, $p\in\{0.3,0.5,0.7\}$, base frequency, fade curve, and base noise type. \textbf{Measured:} (i)~GPU frame time (OpenGL timer queries, fixed resolution, averaged over many frames); (ii)~RMSE of $N$ octaves against a 12-octave reference, plus side-by-side renderings; (iii)~the radially and angularly averaged power spectrum of rendered fBm images (a small C++ DFT tool on exported images).

\textbf{Expected:} cost grows roughly linearly with $N$; spectral falloff follows a power law set by $p$; the image converges after a few octaves, and octaves above the pixel Nyquist limit add aliasing rather than detail. Value noise should look blocky and grid-aligned, Perlin noise smoother but still somewhat axis-aligned, and simplex noise the most isotropic and cheaper per sample in higher dimensions. Cubic fading should show faint grid-aligned artifacts under bump mapping that quintic removes.

\section{Implementation Plan}\label{sec:plan}
\textbf{Stack:} C++, OpenGL 3.3 core, GLSL, GLFW, GLAD, GLM. \textbf{Input:} no external data; simple meshes (plane, sphere, cube, terrain grid) with positions, UVs and normals; all textures are computed in shaders.
\textbf{Stages:} (1)~scene framework with camera and uniform-driven parameters; (2)~hash-based gradient and value noise behind a common interface, validated by visualizing raw noise; (3)~fBm, turbulence, materials and terrain, plus 2D simplex noise; (4)~bump mapping and animation; (5)~experiment harness (parameter sweeps, timer queries, image export, C++ analysis tools); (6)~report and demo scenes.
\textbf{Simplifications:} an integer hash replaces Perlin's permutation table; we run in real time on the GPU rather than in an offline shading language; octave fading by pixel footprint only if time permits.

\section{Expected Challenges}\label{sec:challenges}
(1)~\emph{Aliasing:} high octaves alias under minification, so a justified per-pixel octave cutoff (e.g.\ from screen-space derivatives) needs care. (2)~\emph{Hash quality and fair measurement:} weak GPU hashes cause repetition or directional artifacts, and timer-query results are noisy, so runs must be controlled and repeated.

\section{Team Responsibilities}\label{sec:team}
\textbf{Keshav:} noise core (hashing, value/Perlin/simplex, interpolants), fBm/turbulence, materials, terrain height field. \textbf{Karanveer:} rendering framework, terrain mesh and shading, bump-mapping/animation (and domain warping if chosen), experiment harness and analysis tools. Both will review each other's code and be able to explain every component.

\begin{thebibliography}{9}
\bibitem{perlin1985} Ken Perlin. 1985. An Image Synthesizer. \emph{Computer Graphics (SIGGRAPH '85)} 19, 3, 287--296. https://doi.org/10.1145/325334.325247
\bibitem{perlin2002} Ken Perlin. 2002. Improving Noise. \emph{ACM Trans.\ Graph. (SIGGRAPH '02)} 21, 3, 681--682.
\bibitem{iq} Inigo Quilez. Articles on procedural graphics. https://iquilezles.org/articles/
\end{thebibliography}

\end{document}
